\documentclass[12pt]{article}
\usepackage[utf8]{inputenc}
\usepackage{amsmath}
\usepackage{graphicx}
\usepackage{natbib}
\usepackage{hyperref}
\usepackage{lipsum} % For dummy text, you can remove this in your final document

\title{Advancements in Renewable Energy Technologies: Solar, Wind, and Tidal Energy}
\author{Your Name}
\date{\today}

\begin{document}

\maketitle

\begin{abstract}
Renewable energy technologies such as solar, wind, and tidal power play a crucial role in transitioning towards sustainable energy systems. This paper reviews current advancements in these technologies, discusses efficiency enhancements, scalability challenges, environmental impacts, and explores future trends in renewable energy innovation.
\end{abstract}

\section{Introduction}
Renewable energy sources offer significant potential for reducing greenhouse gas emissions and mitigating climate change impacts. Solar, wind, and tidal energy technologies are key pillars in the global transition towards sustainable energy systems. This paper explores recent advancements in these renewable energy technologies, focusing on efficiency improvements, scalability challenges, and environmental considerations.

\section{Technology Overview}
\subsection{Solar Energy}
Overview of photovoltaic (PV) technology and concentrating solar power (CSP) systems, including efficiency improvements in solar cell technologies and storage solutions \citep{green2018solar, poluzzi2018concentrating}.

\subsection{Wind Energy}
Advancements in wind turbine design, offshore wind farms, and grid integration strategies to enhance efficiency and reliability \citep{barthelmie2010future, chaudhry2018wind}.

\subsection{Tidal Energy}
Current tidal energy extraction technologies, such as tidal barrages and tidal turbines, addressing challenges in efficiency, environmental impact, and marine ecosystem interactions \citep{bahaj2007tidal, johnson2018tidal}.

\section{Efficiency Enhancements}
\subsection{Solar PV Efficiency}
Improvements in solar cell efficiency through material advancements, such as perovskite solar cells, and innovations in panel design and manufacturing processes \citep{yang2017perovskite, zhao2017review}.

\subsection{Wind Turbine Technology}
Innovations in blade design, rotor configurations, and control systems to optimize energy capture and reduce maintenance costs \citep{wood2011advances, saito2013review}.

\subsection{Tidal Stream Generators}
Enhancements in tidal turbine efficiency, including developments in turbine blade materials, hydrodynamic design, and site-specific optimizations \citep{bahaj2019tidal, lewis2014tidal}.

\section{Scalability}
\subsection{Solar Power}
Challenges and opportunities in scaling up solar PV installations globally, including land use requirements, grid integration issues, and economic feasibility \citep{denholm2009land, feldman2016scaling}.

\subsection{Wind Farms}
Scalability challenges in offshore wind farms, addressing logistics, transmission infrastructure, and environmental impact assessments \citep{musial2010assessment, zhou2019towards}.

\subsection{Tidal Energy Arrays}
Feasibility and environmental implications of scaling up tidal energy arrays, optimizing array layouts, and addressing potential impacts on marine ecosystems \citep{bahaj2013impact, neill2018tidal}.

\section{Future Trends}
\subsection{Technological Innovations}
Emerging technologies such as floating solar farms, next-generation wind turbine designs, and underwater tidal energy platforms \citep{chabot2018floating, musslewhite2016underwater}.

\subsection{Integration with Energy Storage}
Advancements in energy storage technologies, including batteries and hydrogen storage, to enhance renewable energy reliability and grid stability \citep{noel2018overview, zakeri2015review}.

\section*{References}
\bibliographystyle{plainnat}
\bibliography{prompt15_35}

\end{document}
